\begin{abstract}
\noindent Estimates of the reproduction number from surveillance counts are usually assumed to become more precise as more cases are reported, although transmission varies between periods. Using a branching process with environmental variation in transmission, negative binomial offspring and binomial reporting, we show that the error variance of the mean reproduction number per period is $V(m)=v_R+A/m$ for $m$ reported cases. Precision gains therefore diminish beyond the balance scale $\mx=A/v_R$, and the budget-optimal case volume is $m^*=\sqrt{(c_0/c_1)\mx}$ per period for a cost ratio $c_0/c_1$ of periods to cases. We analysed US influenza hospital admissions (10506 state-weeks, 2020--2026). In the early-growth windows of four seasons, the relative loss of a renewal-type ratio declined with admissions towards a positive intercept. This gave a balance scale of 37.4 admissions per week on the scale of that ratio (95\% CI: 19.1, 71.0). A two-step inversion through data-matched simulations mapped this ratio and a weighted growth ratio to latent environmental variances of 0.040 (95\% CI: 0.026, 0.056) and 0.043 (95\% CI: 0.025, 0.064). The two estimates were consistent (ratio 0.93). Both are conditional on a 3.6-day generation interval and on independent weekly shocks, as the lag-1 autocorrelation of the data supported. Of the intercept, 41\% to 51\% was common to all states in a calendar week. Under mandatory reporting, facility completeness changed the intercept by 2\% (upper bound 5.3\%; equivalence margin 17\% at reference volatility 0.10). Weekly growth was 1.58 times (95\% CI: 0.95, 2.23) as dispersed in the 2024 voluntary period as in the same months of 2023. Where weekly admissions exceed the balance scale, as in 55\% of state-weeks in the growth windows, further periods and regions improve precision more than further cases per period. High facility participation was associated with lower volatility.

\medskip
\noindent\textbf{Keywords:}\\
Reproduction number\\
Environmental stochasticity\\
Surveillance design\\
Reporting heterogeneity\\
Influenza hospitalisation\\
Diminishing returns
\end{abstract}
